\documentclass[11pt]{beamer}

\usetheme{Madrid}
\usecolortheme{default}
\definecolor{unicaucaverde}{RGB}{0,90,60}
\setbeamercolor{structure}{fg=unicaucaverde}
\setbeamertemplate{navigation symbols}{}

\usepackage[utf8]{inputenc}
\usepackage[spanish]{babel}
\usepackage{amsmath,amssymb}
\usepackage{listings}
\usepackage{xcolor}

\lstdefinestyle{cstyle}{
  language=C,
  basicstyle=\ttfamily\scriptsize,
  keywordstyle=\color{unicaucaverde}\bfseries,
  commentstyle=\color{gray}\itshape,
  stringstyle=\color{orange!70!black},
  breaklines=true,
  showstringspaces=false,
  tabsize=2
}
\lstset{style=cstyle}

\title[Semana 3: Funciones]{Funciones y modularidad}
\subtitle{Programación --- Semana 3 de 16}
\author{Universidad del Cauca}
\date{2026}

\begin{document}

\begin{frame}
  \titlepage
\end{frame}

\begin{frame}{Semana 3: dónde estamos}
  \begin{itemize}
    \item Continuación de Fundamentos de programación en C.
    \item Subtemas 3.1--3.2 --- RA3.1, RA3.2.
    \item 4 horas: funciones y diseño descendente, luego paso por valor y ámbito.
  \end{itemize}
  \tableofcontents
\end{frame}

\section{Funciones}

\begin{frame}{¿Qué es una función?}
  \begin{itemize}
    \item Bloque de código con nombre, reutilizable.
    \item Puede recibir \textbf{parámetros} y devolver un \textbf{valor de retorno}.
    \item \textbf{Diseño descendente}: dividir un problema grande en subproblemas, cada uno
          resuelto por su propia función.
  \end{itemize}
  \begin{block}{Ventajas}
    Reutilización, legibilidad, mantenibilidad, prueba independiente.
  \end{block}
\end{frame}

\begin{frame}[fragile]{Prototipo y definición}
\begin{lstlisting}
double areaCirculo(double radio);   // prototipo

int main(void) {
    double area = areaCirculo(3.0);
    printf("Area = %.2f\n", area);
    return 0;
}

double areaCirculo(double radio) {  // definicion
    return 3.14159 * radio * radio;
}
\end{lstlisting}
\end{frame}

\begin{frame}[fragile]{Funciones \texttt{void}}
\begin{lstlisting}
void imprimirTabla(int n) {
    for (int i = 1; i <= 10; i++) {
        printf("%d x %d = %d\n", n, i, n * i);
    }
}
\end{lstlisting}
  \begin{itemize}
    \item No devuelve ningún valor: se usa para \emph{hacer algo}, no para \emph{calcular algo}.
    \item Sin \texttt{return valor;} (a lo sumo \texttt{return;} vacío).
  \end{itemize}
  \vfill
  \textbf{Taller (hora 2):} \texttt{areaCirculo}, \texttt{esPar}, \texttt{maximo},
  \texttt{imprimirTabla}.
\end{frame}

\section{Paso por valor y ámbito}

\begin{frame}[fragile]{Paso de parámetros por valor}
\begin{lstlisting}
void incrementar(int x) {
    x = x + 1;   // modifica solo la copia local
}

int main(void) {
    int numero = 5;
    incrementar(numero);
    printf("%d\n", numero);   // imprime 5, NO 6
    return 0;
}
\end{lstlisting}
  \begin{alertblock}{Clave}
    El parámetro recibe una \textbf{copia} del argumento. Modificarlo no afecta la variable
    original del llamador.
  \end{alertblock}
\end{frame}

\begin{frame}[fragile]{Ámbito: variables locales}
  \begin{itemize}
    \item Declaradas \emph{dentro} de una función (incluidos sus parámetros).
    \item Solo existen mientras la función se ejecuta.
    \item Dos funciones pueden tener variables locales con el mismo nombre sin interferir.
  \end{itemize}
\end{frame}

\begin{frame}[fragile]{Ámbito: variables globales}
\begin{lstlisting}
int contadorEventos = 0;   // global

void registrarEvento(void) {
    contadorEventos++;     // misma variable para todas las funciones
}

int main(void) {
    registrarEvento();
    registrarEvento();
    printf("%d\n", contadorEventos);   // imprime 2
    return 0;
}
\end{lstlisting}
  \begin{alertblock}{Cuidado}
    Útiles cuando de verdad hay que compartir estado, pero abusar de ellas dificulta depurar:
    preferir parámetros y valores de retorno.
  \end{alertblock}
  \vfill
  \textbf{Taller (hora 4):} programa modular con variable global compartida.
\end{frame}

\begin{frame}{Cierre de la semana}
  \begin{itemize}
    \item RA3.1: funciones, prototipo/definición, \texttt{void}, diseño descendente.
    \item RA3.2: paso por valor, ámbito local vs. global.
    \item Taller calificado al cierre de la Hora 4.
  \end{itemize}
\end{frame}

\end{document}
